\section{Hardware Support for Parallelism}

\begin{remark}
    \textbf{Software algorithms} such as Bakery or Filter locks have \textbf{linear space complexity} $O(n)$, making them impractical for large numbers of processes. Hardware support is needed to implement locks efficiently.
\end{remark}

\begin{remark}
    Moder architectures offer \textbf{atomic instructions} that allow for operations to happen in a single atomic step. This allows mutex-implementation with \textbf{constant space complexity} $O(1)$ and can be used for lock-free programming.
\end{remark}

\subsection{Typical Hardware Instructions}

\begin{definition}
    \textbf{Compare and Swap (CAS):} Atomically compares the value of a memory location with a given value. If they are equal, it updates the memory location with a new value. Otherwise, it does nothing.
\end{definition}

\begin{definition}
    \textbf{Test and Set (TAS):} Atomically reads the value of a memory location and sets it to 1.
\end{definition}

\begin{definition}
    \textbf{Load Linked / Store Conditional (LL/SC): } LL reads the value of a memory location and sets a "link" to it. SC writes a new value to the memory location only if the link is still valid (i.e., no other process has written to that location since the LL).
\end{definition}

\begin{remark}
    \textbf{Java} offers these funcitonalities via the classes \textbf{AtomicInteger}, \textbf{AtomicLong}, etc.
\end{remark}

\subsection{Locks using atomic operations}

\begin{codeexample}
//TAS LOCK IN Java
public class TASLock implements Lock {
    AtomicBoolean state = new AtomicBoolean(false);

    public void lock() {
        while (state.getAndSet(true)) {
            //busy wait
        }
    }

    public void unlock() {
        state.set(false);
    }
}
\end{codeexample}

\begin{remark}
    This lock consistently checks the state variable with an atomic write operation, which causes a lot of cache coherence traffic. To avoid this we can use TATAS lock (checking if the state is false before trying to acquire the lock).
\end{remark}

\begin{codeexample}
//TAS LOCK IN Java
public class TATASLock implements Lock {
    AtomicBoolean state = new AtomicBoolean(false);

    public void lock() {
        do {
            while (state.get());
        } while (state.getAndSet(true));
    }

    public void unlock() {
        state.set(false);
    }
}
\end{codeexample}

\begin{remark}
    To further optimize this lock we can use a backoff strategy (e.g. exponential backoff) to reduce the cache coherence traffic. This means that if a process fails to acquire the lock, it will wait for a certain amount of time (ex. exponentially increasing) before trying again.
\end{remark}

